% Generated draft: requires mathematical and source review.
\documentclass[11pt]{article}
\usepackage[margin=25mm]{geometry}
\usepackage{fontspec}
\setmainfont{Latin Modern Roman}
\usepackage{amsmath,amssymb,mathtools,unicode-math}
\setmathfont{Latin Modern Math}
\usepackage{enumitem,needspace,xurl,hyperref}
\setlength{\parindent}{0pt}
\setlength{\parskip}{5pt}
\setlength{\emergencystretch}{3em}
\newcommand{\N}{\mathbb{N}}
\newcommand{\Z}{\mathbb{Z}}
\newcommand{\Q}{\mathbb{Q}}
\newcommand{\R}{\mathbb{R}}
\newcommand{\C}{\mathbb{C}}
\newcommand{\F}{\mathbb{F}}
\newcommand{\A}{\mathbb{A}}
\newcommand{\PP}{\mathbb{P}}
\newcommand{\E}{\mathbb{E}}
\newcommand{\eps}{\varepsilon}
\newcommand{\dd}{\,\mathrm{d}}
\begin{document}
\section*{UnsolvedMath ID: 30001443}
\textbf{Source problem code:} catalogue entry
\subsection{Definitions and mathematical statement}
For a complex \(C^*\)-algebra \(A\), let \(\operatorname{Prim}(A)\) denote the set of kernels of its nonzero irreducible \(*\)-representations, equipped with the hull--kernel topology: its closed sets are the sets
\[
\operatorname{hull}(I)=\{P\in\operatorname{Prim}(A):I\subseteq P\},
\]
where \(I\) ranges over the closed two-sided ideals of \(A\). Here amenability of a \(C^*\)-algebra is understood in its usual sense, equivalently nuclearity. Explicitly, nuclearity means that for every finite subset \(F\subseteq A\) and every \(\varepsilon>0\), there exist an integer \(n\geq1\) and completely positive contractive maps
\[
A\xrightarrow{\phi}M_n(\C)\xrightarrow{\psi}A
\quad\text{such that}\quad
\|\psi\phi(a)-a\|<\varepsilon\qquad(a\in F).
\]
A map is completely positive if all its matrix amplifications preserve positive elements.

A topological space is sober if every nonempty irreducible closed subset is the closure of exactly one point. A nonempty closed subset is irreducible if it is not the union of two proper closed subsets. All subsets below carry their subspace topologies, and \(\cong\) denotes homeomorphism.

\textbf{Supplied extension question.} For every locally compact sober topological space \(X\) and every open subset \(U\subseteq X\), is the following implication true?
\[
\begin{gathered}
\exists\text{ amenable }C^*\text{-algebras }A,B:\quad
U\cong\operatorname{Prim}(A),\qquad
X\setminus U\cong\operatorname{Prim}(B)
\\[2pt]
\Longrightarrow\quad
\exists\text{ an amenable }C^*\text{-algebra }C:\quad
X\cong\operatorname{Prim}(C).
\end{gathered}
\]
No separability, unitality, nonzeroness, or Hausdorff hypothesis is supplied, and no restriction excluding \(U=\varnothing\) or \(U=X\) is stated. The conclusion asks only for a realization of \(X\); it does not require an algebra extension involving the particular witnesses \(A\) and \(B\).

\textbf{Unresolved qualification.} The source does not specify its convention for local compactness in a possibly non-Hausdorff space. For example, requiring each point to have a compact neighborhood and requiring a neighborhood base of compact neighborhoods need not express the same condition in this setting. The phrase ``locally compact'' is retained pending verification of the intended convention; no additional separation axiom is imposed.
\subsection{Short English statement}
If an open part of a locally compact sober space and its closed complement can each be realized as the primitive ideal space of an amenable \(C^*\)-algebra, must the whole space admit such a realization? The supplied record labels this question open. The intended non-Hausdorff convention for local compactness requires review.
\subsection{Sources}
\begin{itemize}
\item Local UnsolvedMath registry: supplied complete source record, permanent ID \(30001443\), title ``Gluing Primitive Ideal Spaces of Amenable Algebras,'' exact source problem code \texttt{catalogue entry}, in the Oberwolfach Reports Open Problems collection. The record gives proposed year 2010 and names no individual proposer.
\item Supplied report citation: \emph{C*-Algebren} (2010), Oberwolfach Reports, pp.~667--732, DOI \texttt{10.4171/owr/2010/13}. \url{https://doi.org/10.4171/owr/2010/13}. The background describes this collection as extracted from workshop reports and their open-problem discussions; no problem-specific page is supplied.
\item Supplied literature source, labeled ``original OWR report'': \url{https://doi.org/10.4171/owr/2010/13}. Literature assessment dated \texttt{2026-08-21}: the supplied curation describes the extension conjecture as an open route in the report and states that no authoritative later proof or counterexample to the general formulation was found. This is supplied dated curation, not an independently verified literature assessment; the record requests independent verification of time-sensitive status and cited claims.
\item The registry labels the statement \texttt{corrected\_verified}, explaining that math delimiters were added and a trailing extraction artifact was removed without changing mathematical content. This supplied editorial label does not establish independent mathematical or source verification of this draft.
\end{itemize}
\end{document}
